\IfFileExists{stacks-project.cls}{%
\documentclass{stacks-project}
}{%
\documentclass{amsart}
}

% For dealing with references we use the comment environment
\usepackage{verbatim}
\newenvironment{reference}{\comment}{\endcomment}
%\newenvironment{reference}{}{}
\newenvironment{slogan}{\comment}{\endcomment}
\newenvironment{history}{\comment}{\endcomment}

% For commutative diagrams we use Xy-pic
\usepackage[all]{xy}

% We use 2cell for 2-commutative diagrams.
\xyoption{2cell}
\UseAllTwocells

% We use multicol for the list of chapters between chapters
\usepackage{multicol}

% This is generally recommended for better output
\usepackage{lmodern}
\usepackage[T1]{fontenc}

% For cross-file-references
\usepackage{xr-hyper}

% Package for hypertext links:
\usepackage{hyperref}

% For any local file, say "hello.tex" you want to link to please
% use \externaldocument[hello-]{hello}
\externaldocument[introduction-]{introduction}
\externaldocument[conventions-]{conventions}
\externaldocument[sets-]{sets}
\externaldocument[categories-]{categories}
\externaldocument[topology-]{topology}
\externaldocument[sheaves-]{sheaves}
\externaldocument[sites-]{sites}
\externaldocument[stacks-]{stacks}
\externaldocument[fields-]{fields}
\externaldocument[algebra-]{algebra}
\externaldocument[brauer-]{brauer}
\externaldocument[homology-]{homology}
\externaldocument[derived-]{derived}
\externaldocument[simplicial-]{simplicial}
\externaldocument[more-algebra-]{more-algebra}
\externaldocument[smoothing-]{smoothing}
\externaldocument[modules-]{modules}
\externaldocument[sites-modules-]{sites-modules}
\externaldocument[injectives-]{injectives}
\externaldocument[cohomology-]{cohomology}
\externaldocument[sites-cohomology-]{sites-cohomology}
\externaldocument[dga-]{dga}
\externaldocument[dpa-]{dpa}
\externaldocument[sdga-]{sdga}
\externaldocument[hypercovering-]{hypercovering}
\externaldocument[schemes-]{schemes}
\externaldocument[constructions-]{constructions}
\externaldocument[properties-]{properties}
\externaldocument[morphisms-]{morphisms}
\externaldocument[coherent-]{coherent}
\externaldocument[divisors-]{divisors}
\externaldocument[limits-]{limits}
\externaldocument[varieties-]{varieties}
\externaldocument[topologies-]{topologies}
\externaldocument[descent-]{descent}
\externaldocument[perfect-]{perfect}
\externaldocument[more-morphisms-]{more-morphisms}
\externaldocument[flat-]{flat}
\externaldocument[groupoids-]{groupoids}
\externaldocument[more-groupoids-]{more-groupoids}
\externaldocument[etale-]{etale}
\externaldocument[chow-]{chow}
\externaldocument[intersection-]{intersection}
\externaldocument[pic-]{pic}
\externaldocument[weil-]{weil}
\externaldocument[adequate-]{adequate}
\externaldocument[dualizing-]{dualizing}
\externaldocument[duality-]{duality}
\externaldocument[discriminant-]{discriminant}
\externaldocument[derham-]{derham}
\externaldocument[local-cohomology-]{local-cohomology}
\externaldocument[algebraization-]{algebraization}
\externaldocument[curves-]{curves}
\externaldocument[resolve-]{resolve}
\externaldocument[models-]{models}
\externaldocument[functors-]{functors}
\externaldocument[equiv-]{equiv}
\externaldocument[pione-]{pione}
\externaldocument[etale-cohomology-]{etale-cohomology}
\externaldocument[proetale-]{proetale}
\externaldocument[relative-cycles-]{relative-cycles}
\externaldocument[more-etale-]{more-etale}
\externaldocument[trace-]{trace}
\externaldocument[crystalline-]{crystalline}
\externaldocument[spaces-]{spaces}
\externaldocument[spaces-properties-]{spaces-properties}
\externaldocument[spaces-morphisms-]{spaces-morphisms}
\externaldocument[decent-spaces-]{decent-spaces}
\externaldocument[spaces-cohomology-]{spaces-cohomology}
\externaldocument[spaces-limits-]{spaces-limits}
\externaldocument[spaces-divisors-]{spaces-divisors}
\externaldocument[spaces-over-fields-]{spaces-over-fields}
\externaldocument[spaces-topologies-]{spaces-topologies}
\externaldocument[spaces-descent-]{spaces-descent}
\externaldocument[spaces-perfect-]{spaces-perfect}
\externaldocument[spaces-more-morphisms-]{spaces-more-morphisms}
\externaldocument[spaces-flat-]{spaces-flat}
\externaldocument[spaces-groupoids-]{spaces-groupoids}
\externaldocument[spaces-more-groupoids-]{spaces-more-groupoids}
\externaldocument[bootstrap-]{bootstrap}
\externaldocument[spaces-pushouts-]{spaces-pushouts}
\externaldocument[spaces-chow-]{spaces-chow}
\externaldocument[groupoids-quotients-]{groupoids-quotients}
\externaldocument[spaces-more-cohomology-]{spaces-more-cohomology}
\externaldocument[spaces-simplicial-]{spaces-simplicial}
\externaldocument[spaces-duality-]{spaces-duality}
\externaldocument[formal-spaces-]{formal-spaces}
\externaldocument[restricted-]{restricted}
\externaldocument[spaces-resolve-]{spaces-resolve}
\externaldocument[formal-defos-]{formal-defos}
\externaldocument[defos-]{defos}
\externaldocument[cotangent-]{cotangent}
\externaldocument[examples-defos-]{examples-defos}
\externaldocument[algebraic-]{algebraic}
\externaldocument[examples-stacks-]{examples-stacks}
\externaldocument[stacks-sheaves-]{stacks-sheaves}
\externaldocument[criteria-]{criteria}
\externaldocument[artin-]{artin}
\externaldocument[quot-]{quot}
\externaldocument[stacks-properties-]{stacks-properties}
\externaldocument[stacks-morphisms-]{stacks-morphisms}
\externaldocument[stacks-limits-]{stacks-limits}
\externaldocument[stacks-cohomology-]{stacks-cohomology}
\externaldocument[stacks-perfect-]{stacks-perfect}
\externaldocument[stacks-introduction-]{stacks-introduction}
\externaldocument[stacks-more-morphisms-]{stacks-more-morphisms}
\externaldocument[stacks-geometry-]{stacks-geometry}
\externaldocument[moduli-]{moduli}
\externaldocument[moduli-curves-]{moduli-curves}
\externaldocument[examples-]{examples}
\externaldocument[exercises-]{exercises}
\externaldocument[guide-]{guide}
\externaldocument[desirables-]{desirables}
\externaldocument[coding-]{coding}
\externaldocument[obsolete-]{obsolete}
\externaldocument[fdl-]{fdl}
\externaldocument[index-]{index}

% Theorem environments.
%
\theoremstyle{plain}
\newtheorem{theorem}[subsection]{Theorem}
\newtheorem{proposition}[subsection]{Proposition}
\newtheorem{lemma}[subsection]{Lemma}

\theoremstyle{definition}
\newtheorem{definition}[subsection]{Definition}
\newtheorem{example}[subsection]{Example}
\newtheorem{exercise}[subsection]{Exercise}
\newtheorem{situation}[subsection]{Situation}

\theoremstyle{remark}
\newtheorem{remark}[subsection]{Remark}
\newtheorem{remarks}[subsection]{Remarks}

\numberwithin{equation}{subsection}

% Macros
%
\def\lim{\mathop{\mathrm{lim}}\nolimits}
\def\colim{\mathop{\mathrm{colim}}\nolimits}
\def\Spec{\mathop{\mathrm{Spec}}}
\def\Hom{\mathop{\mathrm{Hom}}\nolimits}
\def\Ext{\mathop{\mathrm{Ext}}\nolimits}
\def\SheafHom{\mathop{\mathcal{H}\!\mathit{om}}\nolimits}
\def\SheafExt{\mathop{\mathcal{E}\!\mathit{xt}}\nolimits}
\def\Sch{\mathit{Sch}}
\def\Mor{\mathop{\mathrm{Mor}}\nolimits}
\def\Ob{\mathop{\mathrm{Ob}}\nolimits}
\def\Sh{\mathop{\mathit{Sh}}\nolimits}
\def\NL{\mathop{N\!L}\nolimits}
\def\CH{\mathop{\mathrm{CH}}\nolimits}
\def\proetale{{pro\text{-}\acute{e}tale}}
\def\etale{{\acute{e}tale}}
\def\QCoh{\mathit{QCoh}}
\def\Ker{\mathop{\mathrm{Ker}}}
\def\Im{\mathop{\mathrm{Im}}}
\def\Coker{\mathop{\mathrm{Coker}}}
\def\Coim{\mathop{\mathrm{Coim}}}

% Boxtimes
%
\DeclareMathSymbol{\boxtimes}{\mathbin}{AMSa}{"02}

%
% Macros for moduli stacks/spaces
%
\def\QCohstack{\mathcal{QC}\!\mathit{oh}}
\def\Cohstack{\mathcal{C}\!\mathit{oh}}
\def\Spacesstack{\mathcal{S}\!\mathit{paces}}
\def\Quotfunctor{\mathrm{Quot}}
\def\Hilbfunctor{\mathrm{Hilb}}
\def\Curvesstack{\mathcal{C}\!\mathit{urves}}
\def\Polarizedstack{\mathcal{P}\!\mathit{olarized}}
\def\Complexesstack{\mathcal{C}\!\mathit{omplexes}}
% \Pic is the operator that assigns to X its picard group, usage \Pic(X)
% \Picardstack_{X/B} denotes the Picard stack of X over B
% \Picardfunctor_{X/B} denotes the Picard functor of X over B
\def\Pic{\mathop{\mathrm{Pic}}\nolimits}
\def\Picardstack{\mathcal{P}\!\mathit{ic}}
\def\Picardfunctor{\mathrm{Pic}}
\def\Deformationcategory{\mathcal{D}\!\mathit{ef}}
\usepackage{cleveref}
\crefname{originalreference}{original reference}{original references}
\Crefname{originalreference}{Original reference}{Original references}
\makeatletter
\newlabel{categories-lemma-fibred-in-groupoids-over-fibred-in-groupoids}{{09WW}{source}{}{sourcepointer-1}{}}
\newlabel{categories-lemma-fibred-in-groupoids-over-fibred-in-groupoids@cref}{{[originalreference][0][]09WW}{[1][0][]source}{}{}{}}
\makeatother
\crefname{originalreference}{original reference}{original references}
\Crefname{originalreference}{Original reference}{Original references}
\makeatletter
\newlabel{categories-lemma-fibred-over-fibred}{{09WV}{source}{}{sourcepointer-2}{}}
\newlabel{categories-lemma-fibred-over-fibred@cref}{{[originalreference][0][]09WV}{[1][0][]source}{}{}{}}
\makeatother
\crefname{originaldefinition}{original definition}{original definitions}
\Crefname{originaldefinition}{Original definition}{Original definitions}
\makeatletter
\newlabel{definition-descent-data}{{026B}{source}{}{sourcepointer-3}{}}
\newlabel{definition-descent-data@cref}{{[originaldefinition][0][]026B}{[1][0][]source}{}{}{}}
\makeatother
\crefname{originaldefinition}{original definition}{original definitions}
\Crefname{originaldefinition}{Original definition}{Original definitions}
\makeatletter
\newlabel{definition-effective-descent-datum}{{026E}{source}{}{sourcepointer-4}{}}
\newlabel{definition-effective-descent-datum@cref}{{[originaldefinition][0][]026E}{[1][0][]source}{}{}{}}
\makeatother
\crefname{originaldefinition}{original definition}{original definitions}
\Crefname{originaldefinition}{Original definition}{Original definitions}
\makeatletter
\newlabel{definition-stack}{{026F}{source}{}{sourcepointer-5}{}}
\newlabel{definition-stack@cref}{{[originaldefinition][0][]026F}{[1][0][]source}{}{}{}}
\makeatother
\crefname{originallemma}{original lemma}{original lemmas}
\Crefname{originallemma}{Original lemma}{Original lemmas}
\makeatletter
\newlabel{lemma-presheaf-mor-map-fibred-categories}{{042V}{source}{}{sourcepointer-6}{}}
\newlabel{lemma-presheaf-mor-map-fibred-categories@cref}{{[originallemma][0][]042V}{[1][0][]source}{}{}{}}
\makeatother
\title{A stack over a stack is a stack over the base}
\author{The Stacks Project Authors}
\date{stacks-project:stacks.tex:local-snapshot-20261012; GNU FDL 1.2 or later}
\begin{document}
\maketitle
\noindent\textit{Editorial note.} No omission markers in-span; proof complete as written.\par\medskip
\noindent\textit{Allowed standard prerequisites (editorial).}\begin{itemize}
\item Fibred over fibred. categories-lemma-fibred-over-fibred (cross-chapter, named). Gives that Y is a fibred category over C, settling stack condition (1).
\item Stack and descent data definitions. definition-stack, definition-descent-data, definition-effective-descent-datum (same chapter, named). Fix the three stack conditions and the shape of the descent data used in condition (3).
\end{itemize}\medskip
\ifdefined\hypertarget\hypertarget{sourcepointer-1}{}\fi
\noindent\textit{Source reference.} Original reference 09WW: focused-seeds/geometry/S9-stk-stack-over-stack/ref-tags.txt, source lines 5--5. Navigation pointer only: the counted proof cites this statement by \textbackslash{}ref as a standing input; the deduction the counted proof has to make is carried out inside the counted span, and no part of that deduction is delegated to this pointer.\par\medskip
\ifdefined\hypertarget\hypertarget{sourcepointer-2}{}\fi
\noindent\textit{Source reference.} Original reference 09WV: focused-seeds/geometry/S9-stk-stack-over-stack/ref-tags.txt, source lines 6--6. Navigation pointer only: the counted proof cites this statement by \textbackslash{}ref as a standing input; the deduction the counted proof has to make is carried out inside the counted span, and no part of that deduction is delegated to this pointer.\par\medskip
\ifdefined\hypertarget\hypertarget{sourcepointer-3}{}\fi
\noindent\textit{Source reference.} Original definition 026B: sources/stacks-stacks/pinned-original-source.tex, source lines 498--498. Navigation pointer only: the counted proof cites this statement by \textbackslash{}ref as a standing input; the deduction the counted proof has to make is carried out inside the counted span, and no part of that deduction is delegated to this pointer.\par\medskip
\ifdefined\hypertarget\hypertarget{sourcepointer-4}{}\fi
\noindent\textit{Source reference.} Original definition 026E: sources/stacks-stacks/pinned-original-source.tex, source lines 646--646. Navigation pointer only: the counted proof cites this statement by \textbackslash{}ref as a standing input; the deduction the counted proof has to make is carried out inside the counted span, and no part of that deduction is delegated to this pointer.\par\medskip
\ifdefined\hypertarget\hypertarget{sourcepointer-5}{}\fi
\noindent\textit{Source reference.} Original definition 026F: sources/stacks-stacks/pinned-original-source.tex, source lines 785--785. Navigation pointer only: the counted proof cites this statement by \textbackslash{}ref as a standing input; the deduction the counted proof has to make is carried out inside the counted span, and no part of that deduction is delegated to this pointer.\par\medskip
\ifdefined\hypertarget\hypertarget{sourcepointer-6}{}\fi
\noindent\textit{Source reference.} Original lemma 042V: sources/stacks-stacks/pinned-original-source.tex, source lines 388--388. Navigation pointer only: the counted proof cites this statement by \textbackslash{}ref as a standing input; the deduction the counted proof has to make is carried out inside the counted span, and no part of that deduction is delegated to this pointer.\par\medskip
\begin{lemma}
\label{lemma-stack-over-stack}
Let $\mathcal{C}$ be a site. Let $p : \mathcal{X} \to \mathcal{C}$
be a stack. Endow $\mathcal{X}$ with the topology inherited from
$\mathcal{C}$ and let $q : \mathcal{Y} \to \mathcal{X}$ be a stack.
Then $\mathcal{Y}$ is a stack over $\mathcal{C}$.
If $p$ and $q$ define stacks in groupoids, then
$\mathcal{Y}$ is a stack in groupoids over $\mathcal{C}$.
\end{lemma}

\begin{proof}
We check the three conditions in Definition \ref{definition-stack}
to prove that $\mathcal{Y}$ is a stack over $\mathcal{C}$.
By Categories, Lemma \ref{categories-lemma-fibred-over-fibred}
we find that $\mathcal{Y}$ is a fibred category over $\mathcal{C}$.
Thus condition (1) holds.

\medskip\noindent
Let $U$ be an object of $\mathcal{C}$ and let $y_1, y_2$ be objects
of $\mathcal{Y}$ over $U$. Denote $x_i = q(y_i)$ in $\mathcal{X}$.
Consider the map of presheaves
$$
q : \mathit{Mor}_{\mathcal{Y}/\mathcal{C}}(y_1, y_2)
\longrightarrow
\mathit{Mor}_{\mathcal{X}/\mathcal{C}}(x_1, x_2)
$$
on $\mathcal{C}/U$, see Lemma \ref{lemma-presheaf-mor-map-fibred-categories}.
Let $\{U_i \to U\}$ be a covering and let $\varphi_i$ be a section
of the presheaf on the left over $U_i$ such that $\varphi_i$ and
$\varphi_j$ restrict to the same section over $U_i \times_U U_j$.
We have to find a morphism $\varphi : y_1 \to y_2$ restricting to $\varphi_i$.
Note that $q(\varphi_i) = \psi|_{U_i}$ for some morphism
$\psi : x_1 \to x_2$ over $U$ because the second presheaf is a sheaf
(by assumption). Let $y_{12} \to y_2$ be the strongly $\mathcal{X}$-cartesian
morphism of $\mathcal{Y}$ lying over $\psi$. Then $\varphi_i$ corresponds
to a morphism $\varphi'_i : y_1|_{U_i} \to y_{12}|_{U_i}$ over $x_1|_{U_i}$.
In other words, $\varphi'_i$ now define local sections of the presheaf
$$
\mathit{Mor}_{\mathcal{Y}/\mathcal{X}}(y_1, y_{12})
$$
over the members of the covering $\{x_1|_{U_i} \to x_1\}$. By assumption these
glue to a unique morphism $y_1 \to y_{12}$ which composed with the given
morphism $y_{12} \to y_2$ produces the desired morphism $y_1 \to y_2$.

\medskip\noindent
Finally, we show that descent data are effective. Let $\{f_i : U_i \to U\}$
be a covering of $\mathcal{C}$ and let $(y_i, \varphi_{ij})$ be a descent
datum relative to this covering (Definition \ref{definition-descent-data}).
Setting $x_i = q(y_i)$ and $\psi_{ij} = q(\varphi_{ij})$
we obtain a descent datum $(x_i, \psi_{ij})$ for the covering in $\mathcal{X}$.
By assumption on $\mathcal{X}$ we may assume $x_i = x|_{U_i}$
and the $\psi_{ij}$ equal to the canonical descent datum
(Definition \ref{definition-effective-descent-datum}).
In this case $\{x|_{U_i} \to x\}$ is a covering and we can view
$(y_i, \varphi_{ij})$ as a descent datum relative to this covering.
By our assumption that $\mathcal{Y}$ is a stack over $\mathcal{C}$
we see that it is effective which finishes the proof of condition (3).

\medskip\noindent
The final assertion follows because $\mathcal{Y}$ is a stack over
$\mathcal{C}$ and is fibred in groupoids by
Categories, Lemma
\ref{categories-lemma-fibred-in-groupoids-over-fibred-in-groupoids}.
\end{proof}

\end{document}
